%% Theory of Computing
%% File: v018a022.tex
%% Publication date: December 29, 2022
%% Authors: Alexander Golovnev and Ishay Haviv
%
%
%
%

%

%

%

%
%
%
%
%
%
%
%
%
%
%
%
%
%
%
%
%
%
%
%
%
%

%

\documentclass{toc}

%
%
%
\tocdetails{%
  title = {The (Generalized) Orthogonality Dimension of (Generalized) Kneser Graphs: Bounds and Applications},
  runningtitle = {The (Generalized) Orthogonality Dimension of (Generalized) Kneser Graphs}, %
%
%
  number_of_pages = {20},
  number_of_bibitems = {47},
  number_of_figures = {0},
  conference_version = {CCC21}, %
    %
    %
%
%
%
  author = {Alexander Golovnev and Ishay Haviv},
    %
    %
    %
    %
    %
  authorlist = {Alexander Golovnev, Ishay Haviv},
    %
    %
    %
    %
    %
    %
    %
%
%
  %
  %
  %
  %
%
  acmclassification = {F.2.2, G.2.2},
    %
    %
  amsclassification = {05C50, 68Q06, 68Q17, 68R10},
    %
    %
  keywords = {orthogonality dimension, minrank, Kneser graphs, hardness of approximation, circuit complexity},
    %
    %
    %
    %
    %
}   %

    %
    %
    %
    %
    %
    %
    %
    %
    %
    %
    %
    %
    %
    %
    %
    %
    %
    %
    %
    %
    %
    %
    %
    %
    %
    %
    %
    %
    %

%
%
\tocdetails{%
volume = 18,      %
number = 22,       %
year = 2022,
%
%
%
%
%
%
%
%
%
%
%
received = {October 1, 2021},
revised = {August 21, 2022},
published = {December 29, 2022},
%
%
%
doi = {10.4086/toc.2022.v018.a022},
       %
}   %

%
%
%
%
%
%

%
%
%

%


\newcommand{\Fset}{\mathbb{F}}
\newcommand{\NP}{\mathsf{NP}}
\newcommand{\UGC}{\mathsf{UGC}}
\newcommand{\UG}{\mathsf{UG}}   %
\newcommand{\od}{\overline{\xi}}
\newcommand{\eps}{\epsilon}
\newcommand{\rank}{\mathop{\mathrm{rank}}}
\newcommand{\minrank}{\mathop{\mathrm{minrk}}}
\newcommand{\vchrom}{{\chi_\mathrm{v}}}
\newcommand{\calW}{{\cal W}}
\newcommand{\calG}{{\cal G}}
\newcommand{\linspan}{\mathop{\mathrm{span}}}
\newcommand\Kneser[3]{K^<(#1,#2,#3)}
\newcommand\mymod{\mbox{mod}}
\newcommand\fish{{\rhd \hskip -0.5 em <}}

\begin{document}

\begin{frontmatter}%

%
%
%

%
\title{title of paper here}  %

%
%

\title{The (Generalized) Orthogonality Dimension of (Generalized) Kneser Graphs: Bounds and Applications\titlefootnote{A conference version of this paper appeared in
the \href{http://doi.org/10.4230/LIPIcs.CCC.2021.8}{Proceedings of the 36th Computational Complexity Conference, 2021 (CCC'21)}~\cite{conf-version}.}}

%
%
%
%
%
%
%
%
%
%
%
%


%
%
%
%

%

\author[golovnev]{Alexander Golovnev\thanks{Partially supported by a Rabin Postdoctoral Fellowship.}}
\author[haviv]{Ishay Haviv\thanks{Partially supported by the Israel Science Foundation (grant No. 1218/20).}}

%
%
%
%

%
%
%
%
%
%
%
%
\begin{abstract}
  The \emph{orthogonality dimension} of a graph $G=(V,E)$ over a field $\Fset$ is the smallest integer $t$ for which there exists an assignment of a vector $u_v \in \Fset^t$ with $\langle u_v,u_v \rangle \neq 0$ to every vertex $v \in V$, such that
%
%
%
%
$\langle u_v, u_{w} \rangle = 0$ whenever $v$ and $w$  %
are adjacent vertices in $G$.
The study of the orthogonality dimension of graphs is motivated by various applications in information theory and in theoretical computer science.
The contribution of the present work is %
twofold.  %

First, we prove that there exists a constant $c$ such that for every sufficiently large
fixed  %
integer $t$, it is $\NP$-hard to 
%
%
distinguish between the cases that the orthogonality dimension over $\R$ of an input graph is at most $t$ and at least $3t/2-c$.
At the heart of the proof lies a geometric result, which might be of independent interest, on a generalization %
of the orthogonality dimension parameter
(where vertices correspond to \emph{subspaces} rather than vectors)
%
%
for the family of \emph{Kneser graphs}. The result is related to a long-standing graph-theoretic conjecture due to Stahl (J.~Combin. Theory--B,~1976).

Second, we study the smallest possible orthogonality dimension over finite fields of the complement of graphs that do not contain certain fixed subgraphs.
In particular, we provide an explicit construction of triangle-free $n$-vertex graphs whose complement has orthogonality dimension over the binary field at most $n^{1-\delta}$ for some constant $\delta >0$.
Our results involve constructions from
%
a family of  %
\emph{generalized Kneser graphs} %
%
%
%
(namely, threshold Kneser graphs)
and are motivated by the rigidity approach to circuit lower bounds.
We use them to answer
questions raised by Codenotti, Pudl{\'{a}}k, and Resta (Theoret.~Comput.~Sci.,~2000), and in particular, to disprove their Odd Alternating Cycle Conjecture over every finite field.
\end{abstract}

\iffalse %
%
%
%
%
%
%
%
%
%
%
%
%
%
%
%
%
%
%
%
%
%

\tocarxivcategory{cs.CC}

\fi %

\end{frontmatter}

%
%
%

\section{Introduction}

A $t$-dimensional \emph{orthogonal representation} of a graph $G=(V,E)$ over a field $\Fset$ is an assignment of a vector $u_v \in \Fset^t$ with $\langle u_v,u_v \rangle \neq 0$ to every vertex $v \in V$, such that $\langle u_v, u_{v'} \rangle = 0$ whenever $v$ and $v'$ are adjacent vertices in $G$. The \emph{orthogonality dimension} of a graph $G$ over $\Fset$, denoted by $\od(G, \Fset)$, is the smallest integer $t$ for which there exists a $t$-dimensional orthogonal representation of $G$ over~$\Fset$.\footnote{Orthogonal representations of graphs
%
have originally been defined by Lov\'asz~\cite{Lovasz79}
as orthogonal representations of the complement
in our sense,  %
%
\ie, Lov\'asz's  %
definition requires vectors associated with \emph{non-adjacent} vertices to be orthogonal.
We have decided to use here the other definition (following, e.\,g.,~\cite{deWolfThesis,BrietBLPS15,BrietZ17}), but one may view the notation $\od(G,\Fset)$ as standing for $\xi(\overline{G},\Fset)$.}
%
Orthogonality dimension  %
is closely related to several other well-studied graph parameters. In particular, for every graph $G$ and every field $\Fset$, $\od(G, \Fset)$ is sandwiched between the clique number and the chromatic number of $G$, that is, $\omega(G) \leq \od(G, \Fset) \leq \chi(G)$.

Orthogonal representations of graphs have been found useful over the years for various applications in information theory and in theoretical computer science.
They were originally introduced over the real field in a seminal
%
paper by  %
Lov\'asz~\cite{Lovasz79}, where they were used to define the influential Lov\'asz $\vartheta$-function. The latter was used in~\cite{Lovasz79} to determine the Shannon capacity, a notoriously difficult information-theoretic graph parameter, of the cycle on five vertices, and in the
%
subsequent decades  %
%
%
it was successfully applied in combinatorics, algorithms, and complexity theory (see, e.\,g.,~\cite{Tardos88,knuth1994sandwich,Feige97,AlonK98}).
%
%
%
The orthogonality dimension of graphs plays an important role in several areas of computational complexity.
Over finite fields, the orthogonality dimension and its
%
%
%
relaxation due to Haemers~\cite{Haemers79, Haemers81}, %
%
%
called \emph{minrank},  %
%
%
%
are related to Valiant's concept of matrix rigidity~\cite{Valiant77} (see, e.\,g.,~\cite{CodenottiPR00,Riis07,GolovnevRW17}).
Over the complex field, the orthogonality dimension was used in a characterization of the quantum communication complexity of promise equality problems~\cite{deWolfThesis,BrietBLPS15,BrietZ17} and in the study of the quantum chromatic number~\cite{CameronMNSW07,ScarpaS12}.
%
Orthogonality dimension   %
%
has also been %
investigated in the contexts of hardness of approximation~\cite{Peeters96,LangbergS08}, integrality gaps for linear programming~\cite{HuTS17,Haviv18topo}, and algorithms based on %
semidefinite  %
programming~\cite{ChlamtacH14,HavivMFCS19}.

%
In the present paper we study %
two aspects of the orthogonality dimension of graphs.
First, we prove an $\NP$-hardness result for approximating the orthogonality dimension of graphs over the real field $\R$.
At the heart of the proof lies a geometric result, which might be of independent interest, on a \emph{generalization}
of the orthogonality dimension parameter for the family of \emph{Kneser graphs}. The result is related to a long-standing graph-theoretic conjecture due to Stahl~\cite{Stahl76}.
The second aspect of the orthogonality dimension
%
%
considered in this %
paper,  %
motivated by the area of circuit complexity, is that of determining the smallest possible orthogonality dimension  over finite fields of the complement of graphs that do not contain certain fixed subgraphs. In this context, we prove %
an upper bound  %
on the minrank
%
%
over finite fields for 
%
%
a family of \emph{generalized Kneser graphs} (namely, threshold Kneser graphs).
The bound is used to settle %
some %
questions raised by Codenotti, Pudl{\'{a}}k, and Resta in~\cite{CodenottiPR00}
and to disprove their Odd Alternating Cycle Conjecture over every finite field.

\subsection{Our contribution}

\subsubsection{The generalized orthogonality dimension of Kneser graphs}

We start by considering the computational hardness of determining the orthogonality dimension of graphs over the real field $\R$.
The challenge of understanding the hardness of this parameter was posed already in the late eighties by Lov\'asz, Saks, and Schrijver~\cite{LovaszSS89} (see also~\cite[Chapter~10]{LovaszBook}), and yet, the problem is far from being well-understood.
It is easy to see that deciding whether an input graph $G$ satisfies $\od(G,\R) \leq t$ can be solved in polynomial
%
%
time  %
for $t \in \{1,2\}$, and Peeters~\cite{Peeters96} has shown that it is $\NP$-hard for $t \geq 3$.
His result is known to imply that for every $t \geq 6$ it is $\NP$-hard to decide whether an input graph $G$ satisfies $\od(G,\R) \leq t$ or $\od(G,\R) \geq \lceil 4t/3 \rceil$ (see~\cite{HavivMFCS19}).
In the %
present paper,  %
we improve on the $4/3$ multiplicative gap and prove the following.
%
%
%
%
%
%
%
%
\begin{theorem}\label{thm:IntroHardness}
There exist constants $c$ and $t_0$ such that for every fixed integer $t \geq t_0$, it is $\NP$-hard to distinguish between the graphs $G$ that satisfy $\od(G,\R) \leq t$ from those satisfying $\od(G,\R) \geq 3t/2-c$.
\end{theorem}
%

It is worth noting that in order to obtain hardness results for the orthogonality dimension,
%
%
it is natural to employ known hardness results regarding the closely related chromatic number of graphs. Indeed, it is easy to verify (see, e.\,g.,~\cite{HavivMFCS19}) that every graph $G$ satisfies
\[\log_3 \chi(G) \leq \od(G,\R) \leq \chi(G),\]
hence hardness of deciding whether an input graph $G$ satisfies $\chi(G) \leq t_1$ or $\chi(G) \geq t_2$ immediately implies the hardness of deciding whether it satisfies $\od(G,\R) \leq t_1$ or $\od(G,\R) \geq \log_3 t_2$. In particular, a result of Dinur, Mossel, and Regev~\cite{DinurMR06} on the hardness of the chromatic number implies that %
%
%
%
deciding whether a given graph $G$ satisfies $\od(G,\R) \leq 3$ or $\od(G,\R) \geq t$
%
%
%
%
%
is $\UG'$-hard, where $\UG'$ refers to some variant of
``Unique Games'' (namely, the `$\fish$-shaped' variant). %

However, if one is interested in standard $\NP$-hardness for the orthogonality dimension, the state of the art for the hardness of the chromatic number does not imply
%
%
%
any $\NP$-hardness  %
results using the above reasoning, despite some remarkable recent progress~\cite{BartoBKO21,WZ20}.
Moreover, most hardness proofs for the chromatic number crucially use the fact that an upper bound on the independence number of a graph implies a strong lower bound on its chromatic number (namely, $\chi(G) \geq \frac{|V(G)|}{\alpha(G)}$), whereas an analogue of such a statement for the orthogonality dimension does not hold in general (see, e.\,g.,~\cite[Proposition~2.2]{HavivMFCS19}).

One technique for proving hardness results for the chromatic number that can be applied %
to  %
the orthogonality dimension is that of Garey and Johnson~\cite{GareyJ76a}, who have related hardness of graph coloring to the \emph{multichromatic numbers of Kneser graphs}. The
%
$k$-chromatic
number of a graph~$G$, denoted by $\chi_k(G)$, is the smallest number of colors needed in order to assign to every vertex of $G$ a set of $k$ colors so that adjacent vertices are assigned to disjoint sets. Notice that $\chi_1(G)$ is simply the standard chromatic number $\chi(G)$. The family of Kneser graphs is defined as follows.
\begin{definition}[Kneser Graphs]\label{def:Kneser}
For integers $d \geq 2s$, the \emph{Kneser graph} $K(d,s)$ is the graph whose vertices are all the $s$-subsets of $[d]=\{1,\ldots,d\}$, where two sets are adjacent if they are disjoint.
\end{definition}
\noindent
Note that the multichromatic numbers can be defined in terms of Kneser graphs, namely, $\chi_k(G)$ is the smallest integer $d$ for which there exists a homomorphism from $G$ to $K(d,k)$.

In the seventies, Stahl~\cite{Stahl76} made the following conjecture.
\begin{conjecture}[Stahl's Conjecture~\cite{Stahl76}]\label{conj:Stahl}
For all integers $k$ and $d \geq 2s$,
\[\chi_k(K(d,s)) = \Big \lceil \frac{k}{s} \Big \rceil \cdot (d-2s)+2k.\]
\end{conjecture}
\noindent
Stahl's conjecture has received significant attention in the literature over
the years. Even very recently, it was related to
%
%
Hedetniemi's well-known, recently disproved, conjecture~\cite{TradifZ19}. %
Nevertheless, more than forty years since it was proposed, Stahl's conjecture
is still open.
It is known that the right-hand side in \expref{Conjecture}{conj:Stahl} forms an upper bound on $\chi_k(K(d,s))$, and that this bound is tight up to an additive constant that depends solely on $s$~\cite{CGJ78,Stahl98}.
The precise statement of the conjecture was confirmed only for a few special cases. This includes the case of $k=1$ proved by Lov\'asz~\cite{LovaszKneser}, the cases of $s \leq 2$, $k \leq s$, $d=2s+1$, and $k$ divisible by $s$ proved by Stahl~\cite{Stahl76,Stahl98}, and the case of $s=3$ and $k=4$ proved by Garey and Johnson~\cite{GareyJ76a} (extended to $s=3$ with any $k$ in~\cite{Stahl98}). The result of~\cite{GareyJ76a} was combined there with a simple reduction to show that for every $t \geq 6$, it is $\NP$-hard 
%
%
to distinguish for a given graph $G$ between the cases $\chi(G) \leq t$ and $\chi(G) \geq 2t-4$.

The recent %
paper~\cite{HavivMFCS19}  %
has suggested to borrow the reduction of~\cite{GareyJ76a} to prove hardness results for the orthogonality dimension.
%
%
This approach requires the following generalization of orthogonal representations of graphs over the reals.

\begin{definition}[Orthogonal Subspace Representation\footnote{Over the complex field, the definition is equivalent to the notion of a projective representation from~\cite[Definition~6.1]{manvcinska2016quantum}.}]\label{def:ortho_k-subspace}
%
%
A $t$-dimensional \emph{orthogonal $k$-subspace representation} of a graph $G=(V,E)$ is an assignment of a subspace $U_v \subseteq \R^t$ with $\dim (U_v)=k$ to every vertex $v \in V$, such that the subspaces $U_v$ and $U_{v'}$ are orthogonal whenever $v$ and $v'$ are adjacent in $G$. For a graph $G$, let $\od_k(G,\R)$ denote the smallest integer $t$ for which there exists a $t$-dimensional orthogonal $k$-subspace representation of $G$.
\end{definition}
\noindent
Note that for $k=1$, \expref{Definition}{def:ortho_k-subspace} coincides with the orthogonality dimension over the reals, and that for every graph $G$ and every $k$ it holds that $\od_k(G,\R) \leq \chi_k(G)$.

A combination of the hardness result of Peeters~\cite{Peeters96} and the reduction of~\cite{GareyJ76a} implies the following.

\begin{proposition}[{\cite[Theorem~1.3]{HavivMFCS19}}]\label{prop:xi3vs4}
For every graph $F$, it is $\NP$-hard to 
%
%
distinguish for an input graph $G$ between the cases $\od(G,\R) \leq \od_3(F,\R)$ and $\od(G,\R) \geq \od_4(F,\R)$.
\end{proposition}
\noindent
With \expref{Proposition}{prop:xi3vs4} in hand, it is of interest to find graphs $F$ with a large gap between $\od_3(F,\R)$ and $\od_4(F,\R)$.
In
the %
light of \expref{Conjecture}{conj:Stahl}, it is natural to consider the generalized orthogonality dimension
%
%
for the family of Kneser graphs.
For $k=1$, it was shown in~\cite{Haviv18topo} that the standard chromatic number and the standard orthogonality dimension over $\R$ coincide on all Kneser graphs. In addition, a result of Bukh and Cox~\cite[Proposition~23]{BukhC18} implies that for every $d \geq 2s$ and every~$k$, it holds that $\od_k(K(d,s),\R) \geq kd/s$. This implies that
%
$\chi_k$ and $\od_k(\cdot,\R)$
%
%
%
%
coincide on $K(d,s)$ whenever $k$ is divisible by $s$.

In this %
article  %
we initiate a systematic study of the generalized orthogonality dimension
%
%
of Kneser graphs, analogously to \expref{Conjecture}{conj:Stahl}.
Let us already mention that the arguments applied in the study of Stahl's conjecture do not seem to extend to our question.
The main reason is that the proofs in~\cite{Stahl76,GareyJ76a,CGJ78,Stahl98} use Hilton--Milner-type theorems to characterize the possible structures of the independent sets induced by generalized colorings of Kneser graphs, whereas in our setting, orthogonal subspace representations do not naturally induce large independent sets and the problem seems to require a more geometric approach.

The first non-trivial case is that of Kneser graphs $K(d,s)$ with $s=2$, for which we show that the generalized orthogonality dimension
%
is equal to the multichromatic number.  %

\begin{theorem}\label{thm:Intro_s=2}
For all integers $k \geq 1$ and $d \geq 4$,
\[\od_{k}(K(d,2),\R) = \Big \lceil \frac{k}{2} \Big \rceil \cdot (d-4)+2k.\]
\end{theorem}

We proceed by considering a general $s \geq 3$ and prove the following lower bound.

\begin{theorem}\label{thm:Intro_general_s}
For %
all %
integers $k \geq s \geq 3$ there exists $c=c(s,k)$ such that for all integers $d \geq 2s$,
\[\od_{k}(K(d,s),\R) \geq \frac{k- \lceil \frac{k+1}{s} \rceil +1}{s-1} \cdot d-c.\]
\end{theorem}
%
%
%
%
\noindent
The proofs of Theorems~\ref{thm:Intro_s=2} and~\ref{thm:Intro_general_s} use an inductive argument on $d$.
Note that for $k = \ell \cdot s -1$ where $\ell$ is an integer, the bound provided by \expref{Theorem}{thm:Intro_general_s} is tight up to the additive constant $c$.
Indeed, in this case we get that there exists a constant $c$ such that for all integers $d \geq 2s$ it holds that
\[\ell \cdot d -c \leq \od_{\ell \cdot s-1}(K(d,s),\R) \leq \chi_{\ell \cdot s-1}(K(d,s)) \leq \ell \cdot d -2,\]
where the last inequality follows from the fact that the right-hand side in \expref{Conjecture}{conj:Stahl} is an upper bound on the left-hand side~\cite{Stahl76}.
Note further that for the special case of $k=4$ and $s=3$, \expref{Theorem}{thm:Intro_general_s} implies that there exists a constant $c$ such that $\od_4(K(d,3),\R) \geq 3d/2-c$ for every sufficiently large integer $d$. This, combined with \expref{Proposition}{prop:xi3vs4} and with the fact that $\od_3(K(d,3),\R) = d$, yields our hardness result \expref{Theorem}{thm:IntroHardness}.
%
Let us emphasize that the proof relies on elementary ideas in the spirit of the hardness proof of~\cite{GareyJ76a} for the chromatic number rather than on the PCP theory.

It will be interesting to figure out if the bounds given in \expref{Theorem}{thm:Intro_general_s} can be tightened to the quantity given in the right-hand side of \expref{Conjecture}{conj:Stahl}, at least up to an additive term independent of $d$, as is known to hold for the multichromatic numbers~\cite{Stahl98}.
%
%
%
%
%
%
In particular, it %
would  %
be nice to decide whether for all integers $d \geq 6$ it holds that $\od_4(K(d,3),\R) = 2d-4$.
A positive answer would imply that for every $t \geq 6$, it is $\NP$-hard to decide whether an input graph $G$ satisfies $\od(G) \leq t$ or $\od(G) \geq 2t-4$. We remark, however, that the approach suggested by \expref{Proposition}{prop:xi3vs4} for the hardness of the orthogonality dimension cannot yield a multiplicative hardness gap larger than $2$, as it is easy to see that every graph $F$ satisfies $\od_4(F,\R) \leq \od(F,\R) + \od_3(F,\R) \leq 2 \cdot \od_3(F,\R)$.

We note that known relations between the tractable $\vartheta$-function and the orthogonality dimension over the reals~\cite{Lovasz79} imply that the latter can be approximated in polynomial time to within a factor of $O(n/\log^2 n)$, where $n$ stands for the number of vertices.
It would be interesting to decide if for %
%
every  %
$\eps >0$, it is $\NP$-hard to approximate the orthogonality dimension to within a factor of $n^{1-\eps}$, as is the case for the chromatic number~\cite{Zuckerman07}.

%
\subsubsection{The orthogonality dimension of threshold Kneser graphs}
%

We next
%
%
study
the orthogonality dimension over finite fields of the complement of graphs that do not contain some fixed subgraphs.
In fact, in this context we consider %
%
a relaxation %
of the orthogonality dimension,
%
%
%
%
called minrank, that was introduced by Haemers in~\cite{Haemers79, Haemers81} and is defined as follows.

\begin{definition}[Minrank]\label{def:minrank}
Let $D=(V,E)$ be a directed graph on the vertex set $V=[n]$ and let $\Fset$ be a field.
We say that an $n$ by $n$ matrix $M$ over $\Fset$ \emph{represents} $D$ if $M_{i,i} \neq 0$ for every $i \in V$, and $M_{i,j}=0$ for every distinct $i,j \in V$ such that $(i,j) \notin E$.
The \emph{minrank} of $D$ over $\Fset$ is defined as
\[{\minrank}_{\Fset}(D) =  \min\{{\rank}_{\Fset}(M)\mid M \mbox{ represents }D\mbox{ over }\Fset\}.\]
The definition is naturally extended to (undirected) graphs by replacing every undirected edge with two oppositely directed edges.
\end{definition}
\noindent
Note that for every graph $G$ and every field $\Fset$, ${\minrank}_{\Fset}(\overline{G}) \leq \od(G,\Fset)$.\footnote{Indeed, given a $t$-dimensional orthogonal representation of an $n$-vertex graph $G$ over a field $\Fset$,
%
%
let $B \in \Fset^{n \times t}$ denote the matrix
whose rows contain the vectors associated with the vertices of $G$. Then, the $n$ by $n$ matrix $B \cdot B^T$ represents $\overline{G}$ and has rank at most $t$ over $\Fset$, hence ${\minrank}_{\Fset}(\overline{G}) \leq t$.}

%
%

We consider here the question of whether there are graphs with no short odd cycles and yet low minrank over finite fields.
%
%
This question is motivated by an attempted superlinear lower bound on the rigidity of explicit matrices
for target rank $\Omega(n)$%
%
%
~\cite{CodenottiPR00}. The \emph{rigidity} of an $n$ by $n$ matrix $M$ over a field $\Fset$ with respect to a given parameter $r$
(the target rank)  %
%
is the smallest number of entries that one has to change in $M$ in order to reduce its rank over $\Fset$ to below $r$.
%
Roughly speaking, it was shown in~\cite{Valiant77,C08} that
%
%
%
for $n$ by $n$ matrices $A$ with rigidity $\omega(n)$
for $r = \eps \cdot n$, where $\eps> 0$ is a constant,
the computation of the linear transformation $x\mapsto Ax$
%
%
%
%
by series-parallel arithmetic circuits requires superlinear size.
While Valiant has shown that almost all matrices are highly rigid,
%
%
no explicit construction of matrices
%
%
of superlinear rigidity for linear target rank is known.  %

In 2000, Codenotti, Pudl{\'{a}}k, and Resta~\cite{CodenottiPR00} proposed the
\emph{Odd Alternating Cycle Conjecture}, stated below.
By an \emph{alternating odd cycle} we refer to a directed graph which forms an odd cycle when the orientation of the edges is ignored, and such that the orientation of the edges alternates with one exception.
\begin{conjecture}[The Odd Alternating Cycle Conjecture~\cite{CodenottiPR00}]\label{conj:alternating}
For every field $\Fset$ there exist $\eps >0$ and an odd integer $\ell$ such that every $n$-vertex directed graph $D$ with ${\minrank}_{\Fset} (D)  \leq \eps \cdot n$ contains an alternating cycle of length $\ell$.
\end{conjecture}

It was proved in~\cite{CodenottiPR00} that \expref{Conjecture}{conj:alternating} implies, if true, that certain explicit $n$ by $n$ circulant matrices have superlinear rigidity,
%
namely, $\Omega( n \cdot \log^\delta n)$ for a constant $\delta>0$.
%
%
%
%
%
%
%
It was further proved in~\cite{CodenottiPR00} that every $n$-vertex (undirected) graph $G$ that is free of cycles of length $4$ satisfies ${\minrank}_{\R}(G) > n/6$, whereas
%
there are $n$-vertex (undirected) triangle-free graphs $G$ satisfying ${\minrank}_{\Fset}(G) \leq O(n^{3/4})$ for every field $\Fset$.
%
It was left open whether
%
\expref{Conjecture}{conj:alternating} %
holds  %
for larger
%
odd
values of $\ell$.
%
%
The conjecture was recently disproved~\cite{Haviv18free} over the reals, %
but remained open for finite fields which are of %
particular %
interest in circuit complexity.
For the orthogonality dimension over the binary field $\Fset_2$, it was shown in~\cite{CodenottiPR00} that there exist triangle-free $n$-vertex graphs $G$ satisfying $\od(\overline{G},\Fset_2) = n/4+2$. It was asked there whether every $n$-vertex graph $G$ satisfying $\od(\overline{G},\Fset_2) \leq n/4+1$ must contain a triangle.

In this article
we prove %
an upper bound %
on the minrank
%
%
of \emph{threshold Kneser graphs}  %
(see \expref{Definition}{def:GenKneser})  %
over finite fields.  %
In these graphs, the vertices are all the $s$-subsets of %
the  %
universe $[d]$, where two sets are adjacent if their intersection size is smaller than some integer $m$. Note that for $m=1$ we get the standard family of Kneser graphs (see \expref{Definition}{def:Kneser}).
In the proof we modify and extend an argument of~\cite{Haviv18free}, which is based on linear spaces of multivariate polynomials, building on
%
previous work by  %
Alon~\cite{AlonUnion98}. For the precise statement, see \expref{Theorem}{thm:minrk_Kneser}.  %
Now we  %
turn to %
describing %
several applications of our bound.

As a first application, we establish an explicit construction of graphs that do not contain short odd cycles and yet have low minrank over every finite field.
\begin{theorem}\label{thm:Intro_Cycles}
For every odd integer $\ell \geq 3$ there exists $\delta = \delta(\ell) >0$ such that for every sufficiently large integer $n$, there exists an $n$-vertex graph $G$ with no odd cycle of length at most $\ell$ such that for every finite field $\Fset$,
\[{\minrank}_{\Fset}(G) \leq n^{1-\delta}.\]
\end{theorem}
\noindent
\expref{Theorem}{thm:Intro_Cycles} immediately implies that the Odd Alternating Cycle Conjecture is false over every finite field, even for undirected graphs.
This rules out the approach suggested in~\cite{CodenottiPR00} for lower bounds on the rigidity of certain circulant matrices and thus falls into the recent line of non-rigidity results based on the polynomial method (see, e.\,g.,~\cite{AlmanW17,dvir2019matrix,DvirL19}).
We note, however, that the general upper bound of~\cite{DvirL19} on the rigidity of $n$ by $n$ circulant matrices
%
%
%
is insufficient to disprove the Alternating Cycle Conjecture
(because in~\cite{DvirL19} the upper bound is $n^{1+\eps}$
%
where $\eps > 0$ goes to zero at an unspecified (slow) rate
as $n\to\infty$,  %
whereas in~\cite{CodenottiPR00} the rigidity is only claimed to be
%
%
$\Omega(n \cdot \log^\delta n)$ for a constant $\delta > 0$).

We next consider the behavior of the orthogonality dimension over the binary field of the complement of triangle-free graphs.
It is relevant to mention here that in the proof of \expref{Theorem}{thm:Intro_Cycles}, the matrices that imply the stated bound on the minrank are symmetric (see \expref{Remark}{remark:symmetric}). For the binary field, this can be combined with a matrix decomposition result due to Lempel~\cite{Lempel75} to obtain the following theorem, which answers a question of~\cite{CodenottiPR00} negatively.
\begin{theorem}\label{thm:IntroTriangleFree}
There exists a constant $\delta >0$ such that for every sufficiently large integer $n$ there exists a triangle-free $n$-vertex graph $G$ such that $\od(\overline{G},\Fset_2) \leq n^{1-\delta}$.
\end{theorem}

The above result can also be stated in terms of nearly orthogonal systems.
For a field $\Fset$, a system of vectors in $\Fset^m$ is said to be \emph{nearly
orthogonal} if %
%
none of the vectors of the system is isotropic (self-orthogonal) %
and any set of three of them contains an orthogonal pair.
For the real field, it was proved by Rosenfeld~\cite{Rosenfeld91} that every nearly orthogonal system in $\R^m$ has size at most $2m$.
\expref{Theorem}{thm:IntroTriangleFree} shows that the situation is quite different over the binary field. Namely, it implies that there exists a constant $\delta>0$ such that for infinitely many integers $m$ there exists a nearly orthogonal system in $\Fset_2^m$ of size at least $m^{1+\delta}$.
To see this, consider the triangle-free $n$-vertex graph $G$ given by \expref{Theorem}{thm:IntroTriangleFree}, put $m = \od(\overline{G},\Fset_2) \leq n^{1-\delta}$, and observe that the triangle-freeness of $G$ implies that the $n$ vectors assigned by an $m$-dimensional orthogonal representation of $\overline{G}$ over $\Fset_2$ form a nearly orthogonal system in $\Fset_2^m$.

We finally mention that our bound on the minrank
%
%
%
of threshold Kneser graphs  %
can be used to obtain graphs with a constant \emph{vector chromatic number} $\vchrom$ (see \expref{Definition}{def:chi_v}) whose complement has a polynomially large minrank over every finite field.
\begin{theorem}\label{thm:Intro_chi_v}
There exists a constant $\delta >0$ such that for infinitely many integers $n$ there exists an $n$-vertex graph $G$ such that $\vchrom(G) \leq 3$ and yet $\minrank_{\Fset}(\overline{G}) \geq n^\delta$ for every finite field $\Fset$.
\end{theorem}
\noindent
The interest in such graphs comes from the semidefinite programming algorithmic approach applied in~\cite{ChlamtacH14} for approximating the minrank.
%
%
As explained in~\cite{Haviv18}, %
%
the existence of such graphs implies  %
a limitation on this approach, which is based on the constant vector chromatic number of the complement of the instances. \expref{Theorem}{thm:Intro_chi_v} improves on~\cite[Theorem~1.3]{Haviv18} where the bound on the minrank is shown only for sufficiently large finite fields.

\subsection{Outline}
The rest of the paper is organized as follows.
In \expref{Section}{sec:GenODofKneser}, we prove our bounds on the generalized orthogonality dimension
%
%
of Kneser graphs (Theorems~\ref{thm:Intro_s=2} and~\ref{thm:Intro_general_s}) and derive our hardness result (\expref{Theorem}{thm:IntroHardness}).
In \expref{Section}{sec:ODofGenKneser}, we prove our bound on the minrank
%
%
of threshold Kneser graphs  %
over finite fields  %
and deduce Theorems~\ref{thm:Intro_Cycles},~\ref{thm:IntroTriangleFree}, and~\ref{thm:Intro_chi_v}.

\section{The generalized orthogonality dimension of Kneser graphs}\label{sec:GenODofKneser}

In this section we study the generalized orthogonality dimension
%
%
of Kneser graphs, %
\ie,   %
the quantities $\od_k(K(d,s))$ (recall Definitions~\ref{def:Kneser} and~\ref{def:ortho_k-subspace}), and prove Theorems~\ref{thm:Intro_s=2} and~\ref{thm:Intro_general_s}.
We start with a linear algebra lemma that will be useful in our proofs.

\subsection{Linear algebra lemma}

%
%
%
%
\begin{lemma}\label{lem_cor:good_subspace}
For a real vector space $T$, let $U$ be a subspace of $T$ with $\dim(U) = \ell$, let $\calW$ be a finite collection of subspaces of $T$, and let $\ell' \leq \ell$ be an integer satisfying $\dim(U \cap W) \leq \ell'$ for every $W \in \calW$. Then, there exists a subspace~$U'$ of $U$ with $\dim(U')=\ell-\ell'$ such that $\dim(U' \cap W) = 0$ for every $W \in \calW$.
\end{lemma}

%
%

Intuitively, given a subspace $U$ and a collection $\calW$ as in the lemma, a
%
generic  %
subspace $U'$ of $U$ with dimension $\ell-\ell'$ is expected to have a trivial intersection with each of the subspaces of $\calW$, and thus to satisfy the assertion of the lemma.
%
%
For a proof, see, e.\,g.,~\cite[Chapter~3.1.4]{BabaiF20}.
%

%
%
%
%
%
%
%
%
%
%
%
%
%
%
%
%
%
%
%
%
%
%
%
%
%
%
%
%
%
%
%
%
%
%
%
%
%
%
%
%
%
%
%
%
%
%
%
%
%
%
%
%
%
%
%
%
%
%
%
%
%
%
%
%
%

\subsection{The case \texorpdfstring{$s=2$}{s=2}}

%
Now we turn to the proof of %
\expref{Theorem}{thm:Intro_s=2} which determines the generalized orthogonality dimension %
of the  %
Kneser graphs $K(d,s)$ for $s=2$.

\begin{proof}[Proof of \expref{Theorem}{thm:Intro_s=2}]
Fix an integer $k \geq 1$.
For the upper bound, recall that for all integers $d \geq 4$ we have
\[\od_{k}(K(d,2),\R) \leq \chi_k(K(d,2)) = \Big \lceil \frac{k}{2} \Big \rceil \cdot (d-4)+2k,\]
where the equality holds because \expref{Conjecture}{conj:Stahl} is valid for $s=2$~\cite{Stahl98}.

For the lower bound, we consider the induced subgraph of $K(d,2)$, denoted by $K^{-}(d,2)$, obtained from $K(d,2)$ by removing one of its vertices, say, the vertex $\{1,2\}$.
%
Next we prove  %
that for all integers $d \geq 4$ it holds that
\begin{eqnarray}\label{eq:K-,s=2}
\od_{k}(K^{-}(d,2),\R) \geq \Big \lceil \frac{k}{2} \Big \rceil \cdot (d-4)+2k,
\end{eqnarray}
which immediately implies the required lower bound on $\od_{k}(K(d,2),\R)$ as well.
%
We proceed by  %
induction on $d$.
For $d=4$, the graph $K(d,2)$ is a perfect matching on $6$ vertices, hence its subgraph $K^{-}(d,2)$ clearly contains an edge. Since every orthogonal $k$-subspace representation of this graph assigns to the vertices of this edge orthogonal $k$-subspaces it follows that \[\od_k(K^{-}(4,2),\R) \geq 2k,\]
as desired.
Now, fix some $d > 4$. Assuming that~\eqref{eq:K-,s=2} holds for $d-1$,
we %
prove it for $d$.

Recall that the vertex set $V$ of $K^{-}(d,2)$ consists of all the $2$-subsets of $[d]$ except $\{1,2\}$.
Let $(U_A)_{A \in V}$ be a $t$-dimensional orthogonal $k$-subspace representation of $K^{-}(d,2)$.
We proceed by considering the following two cases.

Assume first that there exists some $i \geq 4$ for which
\begin{eqnarray}\label{eq:U_{1,3}}
\dim(U_{\{1,3\}} \cap U_{\{1,i\}}) \geq \Big\lceil \frac{k}{2} \Big\rceil.
\end{eqnarray}
In this case, consider the induced subgraph of $K^{-}(d,2)$ on the vertex set $V'$ obtained from $V$ by removing the vertex $\{3,i\}$ and all the vertices that include the element $1$.
Notice that this subgraph is isomorphic to $K^{-}(d-1,2)$ and that every vertex of $V'$ is disjoint from either $\{1,3\}$ or from $\{1,i\}$ (or both).
This implies that the restriction $(U_A)_{A \in V'}$ of the given assignment to the vertices of $V'$ forms an orthogonal $k$-subspace representation of $K^{-}(d-1,2)$, all of whose subspaces lie in the subspace of $\R^t$ that is orthogonal to $U = U_{\{1,3\}} \cap U_{\{1,i\}}$.
By applying an orthogonal linear transformation from this subspace to $\R^{t-\dim(U)}$, we obtain that
\[ \od_k(K^{-}(d-1,2),\R) \leq t- \dim(U) \leq t- \Big\lceil \frac{k}{2} \Big\rceil,\]
where in the second inequality we have used~\eqref{eq:U_{1,3}}.
Using the induction hypothesis, this implies that
\[ t \geq \od_k(K^{-}(d-1,2),\R) + \Big\lceil \frac{k}{2} \Big\rceil \geq \Big \lceil \frac{k}{2} \Big \rceil \cdot (d-5)+2k + \Big\lceil \frac{k}{2} \Big\rceil = \Big \lceil \frac{k}{2} \Big \rceil \cdot (d-4)+2k,\]
and we are done.

We are left with the case where for every $i \geq 4$ it holds that
\[\dim(U_{\{1,3\}} \cap U_{\{1,i\}}) \leq \Big\lceil \frac{k}{2} \Big\rceil -1.\]
Apply \expref{Lemma}{lem_cor:good_subspace} to the $k$-subspace $U_{\{1,3\}}$ and the collection $\{ U_{\{1,i\}} \mid 4 \leq i \leq d \}$.
It follows that there exists a subspace $U$ of $U_{\{1,3\}}$ with $\dim(U) = k- (\lceil \frac{k}{2} \rceil -1) \geq \lceil \frac{k}{2} \rceil$ such that for every $i \geq 4$ it holds that $\dim(U \cap U_{\{1,i\}}) = 0$.
Consider the induced subgraph of $K^{-}(d,2)$ on the vertex set $V'$ obtained from $V$ by removing the vertex $\{2,3\}$ and all the vertices that include the element $1$. As before, this subgraph is isomorphic to the graph $K^{-}(d-1,2)$.

We define an orthogonal $k$-subspace representation of this graph as follows.
Let $B$ be a set in $V'$.
If $3 \notin B$ we define $\widetilde{U}_B = U_B$.
Otherwise we have $B = \{3,i\}$ for some $i \geq 4$, and we let $\widetilde{U}_{\{3,i\}}$ be the projection of $U_{\{1,i\}}$ to the subspace of $\R^t$ that is orthogonal to $U$.
Note that the fact that $\dim(U \cap U_{\{1,i\}}) = 0$ guarantees that $\dim(\widetilde{U}_{\{3,i\}}) = \dim(U_{\{1,i\}}) = k$.

To prove that the assignment $(\widetilde{U}_B)_{B \in V'}$ forms an orthogonal $k$-subspace representation of the graph, let $B_1$ and $B_2$ be disjoint sets in $V'$.
If $3 \notin B_1 \cup B_2$ then we have $\widetilde{U}_{B_1} = U_{B_1}$ and $\widetilde{U}_{B_2} = U_{B_2}$, so it is clear that $\widetilde{U}_{B_1}$ and $\widetilde{U}_{B_2}$ are orthogonal.
Otherwise, assume without loss of generality that $B_1 = \{3,i\}$ for some $i \geq 4$ and that $3 \notin B_2$.
In this case we have $\widetilde{U}_{B_2} = U_{B_2}$, and since $B_2$ is disjoint from $B_1$ it is also disjoint from $\{1,i\}$ and from $\{1,3\}$, hence $\widetilde{U}_{B_2}$ is orthogonal to both $U_{\{1,i\}}$ and $U_{\{1,3\}}$ as well as to the projection $\widetilde{U}_{B_1}$ of $U_{\{1,i\}}$ to the subspace orthogonal to $U \subseteq U_{\{1,3\}}$.
We get that $\widetilde{U}_{B_1}$ and $\widetilde{U}_{B_2}$ are orthogonal, as required.

Finally, observe that all the subspaces $\widetilde{U}_{B}$ lie in the subspace of $\R^t$ that is orthogonal to $U$.
Indeed, for sets $B$ with $3 \in B$ this follows from the definition of $\widetilde{U}_{B}$, and for the other sets this holds because they are disjoint from $\{1,3\}$.
By applying an orthogonal linear transformation from this subspace to $\R^{t-\dim(U)}$, we obtain that
\[ \od_k(K^{-}(d-1,2),\R) \leq t- \dim(U) \leq t- \Big\lceil \frac{k}{2} \Big\rceil,\]
and as in the previous case, by the induction hypothesis it follows that
$t \geq \lceil \frac{k}{2} \rceil \cdot (d-4)+2k$, completing the proof.
\end{proof}


\subsection{General \texorpdfstring{$s$}{s}}

We now prove \expref{Theorem}{thm:Intro_general_s}
which provides a lower bound on the generalized orthogonality dimension
%
%
of Kneser graphs $K(d,s)$ for $s \geq 3$.
The approach resembles the one applied for $s=2$ in the proof of \expref{Theorem}{thm:Intro_s=2}.

\begin{proof}[Proof of \expref{Theorem}{thm:Intro_general_s}]
Fix integers $k \geq s \geq 3$ and denote $m = \lceil \frac{k+1}{s} \rceil$.
Let $d_0=d_0(s,k)$ be a sufficiently large integer to be determined later.
We apply an induction on $d$.
To do so, we define $c=c(s,k)$ to be sufficiently large, say, $c=\frac{k-m+1}{s-1} \cdot (d_0+s-2)$, so that the statement of the theorem trivially holds for all integers $d \leq d_0+s-2$, and turn to
%
the proof of %
the statement for $d \geq d_0$ assuming that it holds for $d-(s-1)$.

Let $(U_A)_{A \in V}$ be a $t$-dimensional orthogonal $k$-subspace representation of $K(d,s)$.
We start with some notation.
For an $s$-subset $A$ of $[d]$, an element $i \in A$, and an $s$-subset $B$ of $[d]$ satisfying $A \cap B = \{i\}$, we let $\calG_{A,i}(B)$ denote the collection that consists of the set $B$ and all the sets obtained from $B$ by replacing $i$ with some element from $A \setminus \{i\}$. Note that $|\calG_{A,i}(B)|=s$.
We say that a vertex $A$ of $K(d,s)$ is \emph{good} (with respect to the given orthogonal subspace representation) if there exists an $i \in A$ such that for every vertex $B$ satisfying $A \cap B = \{i\}$ it holds that $\dim(U_A \cap U_C) \leq m-1$ for some $C \in \calG_{A,i}(B)$.

Assume first that there exists a good vertex $A$ in $K(d,s)$ associated with an element $i \in A$.
Applying \expref{Lemma}{lem_cor:good_subspace}, we get that there exists a $(k-m+1)$-subspace $U$ of $U_A$ such that for every vertex $B$ satisfying $A \cap B = \{i\}$ it holds that $\dim(U \cap U_C) =0$ for some $C \in \calG_{A,i}(B)$.
We define an orthogonal $k$-subspace representation of the graph $K(d-(s-1),s)$ on the ground set $[d] \setminus (A \setminus \{i\})$ as follows. Let $B$ be an $s$-subset of $[d] \setminus (A \setminus \{i\})$.
If $i \notin B$ we define $\widetilde{U}_B = U_B$. Otherwise, we have $A \cap B = \{i\}$, and we let $\widetilde{U}_B$ be the projection of $U_C$ to the subspace of $\R^t$ orthogonal to $U$, where $C \in \calG_{A,i}(B)$ is a set satisfying $\dim(U \cap U_C) =0$. Note that this condition guarantees that $\dim(\widetilde{U}_B) = \dim(U_C) = k$.

We claim that the subspaces $\widetilde{U}_B$ form an orthogonal $k$-subspace representation of the graph $K(d-(s-1),s)$.
To see this, let $B_1$ and $B_2$ be disjoint $s$-subsets of $[d] \setminus (A \setminus \{i\})$.
If $i \notin B_1 \cup B_2$ then we have $\widetilde{U}_{B_1} = U_{B_1}$ and $\widetilde{U}_{B_2} = U_{B_2}$, so it is clear that $\widetilde{U}_{B_1}$ and $\widetilde{U}_{B_2}$ are orthogonal. Otherwise, assume without loss of generality that $i \in B_1$ and $i \notin B_2$. In this case, $\widetilde{U}_{B_2} = U_{B_2}$, and $\widetilde{U}_{B_1}$ is the projection of $U_C$ to the subspace of $\R^t$ orthogonal to $U$ for some $C \in \calG_{A,i}(B_1)$.
Since $B_2$ is disjoint from $A$, it follows that the subspace $\widetilde{U}_{B_2}$ is orthogonal to $U_A$ as well as to its subspace $U$.
It also follows that $B_2$ is disjoint from every set in $\calG_{A,i}(B_1)$, hence the subspace $\widetilde{U}_{B_2}$ is orthogonal to $U_C$. We get that $\widetilde{U}_{B_2}$ is orthogonal to $\widetilde{U}_{B_1}$, as required.

Now, observe that the above orthogonal $k$-subspace representation of $K(d-(s-1),s)$ lies in the subspace of $\R^t$ that is orthogonal to the $(k-m+1)$-subspace $U$.
Indeed, for sets $B$ with $i \in B$ this follows from the definition of $\widetilde{U}_{B}$, and for the other sets this holds because they are disjoint from $A$.
By applying an orthogonal linear transformation from this subspace to $\R^{t-\dim(U)}$, it follows that \[\od_{k}(K(d-(s-1),s),\R) \leq t-\dim(U) = t-(k-m+1).\]
Using the induction hypothesis, this implies that
\[t \geq \frac{k-m+1}{s-1} \cdot (d-(s-1)) -c + (k-m+1)= \frac{k-m+1}{s-1} \cdot d-c,\]
and we are done.

We are left with the case where no vertex of $K(d,s)$ is good, for which we need the following lemma.

\begin{lemma}\label{lemma:bad_vertex}
If a vertex $A$ of $K(d,s)$ is not good then there exists a nonzero vector $u_A \in U_A$ such that the number of vertices $D$ of $K(d,s)$ for which $U_D$ is not orthogonal to $u_A$ is at most
\[\binom{2s-1}{2} \cdot \binom{d-2}{s-2}.\]
\end{lemma}

We first show how \expref{Lemma}{lemma:bad_vertex} completes the proof of the theorem.
Assume that no vertex of $K(d,s)$ is good, and consider the  following process: We start with the entire vertex set of $K(d,s)$, and in every iteration we choose an arbitrary vertex $A$ associated with its nonzero vector $u_A \in U_A$ from \expref{Lemma}{lemma:bad_vertex} and eliminate all vertices whose subspaces are not orthogonal to $u_A$. The nonzero vectors associated with the chosen vertices are clearly pairwise orthogonal, and their number, just like the number of iterations in the process, is at least
\[\frac{\binom{d}{s}}{\binom{2s-1}{2} \cdot \binom{d-2}{s-2}} \geq \frac{k-m+1}{s-1} \cdot d-c,\]
where the inequality holds for every $d \geq d_0$ assuming that $d_0 = d_0(s,k)$ is sufficiently large (because the left-hand side of the inequality is quadratic in $d$ whereas the right-hand side is linear in $d$).
However, the size of the obtained orthogonal set cannot exceed the dimension $t$, hence
\[t \geq \frac{k-m+1}{s-1} \cdot d-c,\]
and we are done. It remains to prove \expref{Lemma}{lemma:bad_vertex}.

\begin{proof}[Proof of \expref{Lemma}{lemma:bad_vertex}]
Assume that $A$ is not a good vertex of $K(d,s)$ and fix an arbitrary $i \in A$. Then there exists a vertex $B$ satisfying $A \cap B = \{i\}$ such that $\dim(U_A \cap U_C) \geq m$ for every vertex $C \in \calG_{A,i}(B)$.
Denote $\calG_{A,i}(B) = \{C_1,\ldots,C_s\}$, and recall that every set $C_j$ intersects $A$ at one distinct element. For every $j \in [s]$ define $V_j = U_A \cap U_{C_j}$ and $W_j = V_1 + \cdots + V_j$. Note that
\[W_1 \subseteq W_2 \subseteq \cdots \subseteq W_s \subseteq U_A.\]
Since $\dim(W_1) = \dim(V_1) \geq m$ and $\dim(U_A)= k < m \cdot s$, there must exist some $j \in [s-1]$ for which
\begin{eqnarray}\label{eq:<ell}
\dim(W_{j+1}) - \dim(W_j) < m.
\end{eqnarray}
For this $j$, we have
\begin{eqnarray*}
\dim(V_{j+1} \cap W_j) &=& \dim(V_{j+1})+ \dim(W_j)- \dim(V_{j+1}+W_j)
\\ &=& \dim(V_{j+1})+ \dim(W_j) - \dim(W_{j+1}) > m-m=0,
\end{eqnarray*}
where the inequality follows by combining~\eqref{eq:<ell} with the fact that $\dim(V_{j+1}) \geq m$.
This implies that there exists a nonzero vector $u_A$ in $V_{j+1} \cap W_j$.
Observe that
\[u_A \in U_A \cap U_{C_{j+1}} \cap (U_{C_1} +\cdots+U_{C_j}).\]
Now, consider a vertex $D$ of $K(d,s)$ whose subspace $U_D$ is not orthogonal to $u_A$.
It follows that $U_D$ is not orthogonal to $U_A$, to $U_{C_{j+1}}$, and to $U_{C_1} +\cdots+U_{C_j}$, hence $D$ intersects the sets $A$ and $C_{j+1}$ as well as at least one of the sets $C_1,\ldots,C_{j}$.
We claim that $D$ must include at least two elements from $A \cup B$.
Indeed, $D$ intersects $A$ but if $D$ includes from $A \cup B$ only one element and this element belongs to $A$ then $D$ either does not intersect $C_{j+1}$ or does not intersect any of $C_1,\ldots,C_{j}$.
It follows that the number of vertices $D$ for which $U_D$ is not orthogonal to $u_A$ is bounded from above by the number of $s$-subsets of $[d]$ that include at least two elements from the $2s-1$ elements of $A \cup B$. The latter is at most $\binom{2s-1}{2} \cdot \binom{d-2}{s-2}$, as required.
\end{proof}
The proof
of \expref{Theorem}{thm:Intro_general_s}  %
is completed.
\end{proof}

As immediate corollaries of \expref{Theorem}{thm:Intro_general_s}, we obtain the following.
\begin{corollary}\label{cor:ls-1}
For every integers $s \geq 3$ and $\ell \geq 2$ there exists $c=c(s,\ell)$ such that for all integers $d \geq 2s$,
\[\od_{\ell \cdot s-1}(K(d,s),\R) \geq \ell \cdot d-c.\]
\end{corollary}
\noindent
As mentioned before, the bound given in \expref{Corollary}{cor:ls-1} is tight up to the additive constant $c$.

\begin{corollary}\label{cor:3/2}
There exists an absolute constant $c$ such that for all integers $d \geq 6$,
\[\od_{4}(K(d,3),\R) \geq 3d/2-c.\]
\end{corollary}

Equipped with \expref{Corollary}{cor:3/2}, we are ready to deduce \expref{Theorem}{thm:IntroHardness}.

\begin{proof}[Proof of \expref{Theorem}{thm:IntroHardness}]
Let $t$ be a sufficiently large integer. Recall that a result of~\cite{BukhC18} implies that $\od_3(K(t,3),\R) = t$, whereas \expref{Corollary}{cor:3/2} implies that $\od_{4}(K(t,3),\R) \geq 3t/2-c$ for an absolute constant $c$.
Applying \expref{Proposition}{prop:xi3vs4} with $F = K(t,3)$, it follows that it is $\NP$-hard to 
%
%
distinguish between the graphs $G$ that satisfy $\od(G,\R) \leq t$ from those satisfying $\od(G,\R) \geq 3t/2-c$, as desired.
\end{proof}


%
\section{The minrank of threshold Kneser graphs}\label{sec:ODofGenKneser}

In this section we consider a generalization of the family of
Kneser graphs, defined as follows.  %

%
\begin{definition}[Threshold Kneser Graphs]\label{def:GenKneser}
%
  For integers $m \leq s \leq d$, the
%
\emph{threshold Kneser graph}   %
$\Kneser{d}{s}{m}$ is the graph whose vertices are all the $s$-subsets of $[d]$,
%
and   %
two sets, $A,B$, %
are adjacent if $|A \cap B| < m$.
\end{definition}

For this family of graphs, we prove the following upper bound on the minrank
%
%
over finite fields (recall \expref{Definition}{def:minrank}).

\begin{theorem}\label{thm:minrk_Kneser}
For all integers $m \leq s \leq d$ and for every finite field $\Fset$,
\[{\minrank}_{\Fset}(\Kneser{d}{s}{m}) \leq \sum_{i=0}^{s-m}\binom{d}{i}.\]
Moreover, the bound on the minrank can be achieved by a symmetric matrix.
\end{theorem}

\begin{remark}\label{remark:symmetric}
\expref{Theorem}{thm:minrk_Kneser} guarantees that the bound on the minrank can be achieved by a symmetric matrix.
This will be crucial for one of our applications, namely, for a construction of triangle-free graphs whose complement has low orthogonality dimension over the binary field $\Fset_2$ (see \expref{Section}{sec:OD_F_2_triangle}).
We remark, however, that for undirected graphs and for fields of characteristic different from $2$, attaining the bound on the minrank by a symmetric matrix can be achieved easily with a factor of $2$ worse bound on the minrank. Indeed, if a matrix $M$ represents a graph $G$ over a field $\Fset$ of characteristic different from $2$ and satisfies ${\rank}_{\Fset}(M) = r$ then the matrix $M+M^T$ also represents $G$ and has rank at most $2r$ over $\Fset$.
This argument does not hold over fields of characteristic $2$, since in this case the diagonal entries of $M+M^T$ are all %
zero.  %
\end{remark}

As in the previous section, we start with a simple linear algebra lemma.

\subsection{Linear algebra lemma}

\begin{lemma}\label{lemma:rankp_R}
  For a graph $G$ on the vertex set $[n]$, let $M \in \Z^{n \times n}$ be an integer matrix such that $M_{i,i}=1$ for every $i \in [n]$, and $M_{i,j} = 0$ for every
pair of  %
distinct non-adjacent vertices $i$ and $j$ in $G$. Then, for every finite field $\Fset$, ${\minrank}_{\Fset}(G) \leq \rank_\R(M)$.
\end{lemma}

%
\begin{proof}
The lemma follows from the fact that for every finite field $\Fset$, it holds that $\rank_{\Fset} (M) \leq \rank_{\R} (M)$.
\end{proof}

%
%
%
%
%
%
%
%
%
%
%
%
%
%
%
%
%
%
%
%
%
%
%
%
%
%
%
%
%
%
%

\subsection{Proof of \expref{Theorem}{thm:minrk_Kneser}}

We are ready to prove \expref{Theorem}{thm:minrk_Kneser}.

\begin{proof}[Proof of \expref{Theorem}{thm:minrk_Kneser}]
Consider the polynomial $q \in \R[x]$ defined by
\[q(x) = {\binom {x-m} {s-m}} = \frac{1}{(s-m)!} \cdot (x-m)(x-(m+1)) \cdots (x-(s-1)).\]
Observe that $q$ is an integer-valued polynomial of degree $s-m$.
%
%
Let $f: \{0,1\}^d \times \{0,1\}^d \rightarrow \R$ be the function defined by
\[ f(x,y) = q  \Big ( \sum_{i=1}^{d}{x_i y_i}\Big )\]
for every $x,y \in \{0,1\}^d$.
Expanding $f$ as a linear combination of monomials, the relation $z^2 = z$ for $z \in \{0,1\}$ implies that one can reduce to $1$ the exponent of each variable occuring in a monomial. It follows that $f$ can be represented as a multilinear polynomial in the $2d$ variables of $x$ and $y$. By combining terms involving the same monomial in the variables of $x$, one can write $f$ as
\[ f(x,y) = \sum_{i=1}^{R}{g_i(x) h_i(y)} \]
for an integer $R$ and functions $g_i, h_i : \{0,1\}^d \rightarrow \R$, $i \in [R]$, such that the $g_i$'s are distinct multilinear monomials of total degree at most $s-m$ in $d$ variables. It follows that $R \leq \sum_{i=0}^{s-m}\binom{d}{i}$.

Now, let $M_1$ and $M_2$ be the $2^d \times R$ matrices whose rows are indexed by $\{0,1\}^d$ and whose columns are indexed by $[R]$, defined by $(M_1)_{x,i} = g_i(x)$ and $(M_2)_{x,i} = h_i(x)$. Then, the rank over $\R$ of the matrix $M = M_1 \cdot M_2^T$ is at most $R$ and for every $x,y \in\{0,1\}^d$ it holds that $M_{x,y} = f(x,y)$. By the definition of $f$ the matrix $M$ is symmetric, and since $q$ is an integer-valued polynomial, all of its entries are integer.

Finally, let $V$ be the vertex set of $\Kneser{d}{s}{m}$, that is, the collection of all $s$-subsets of $[d]$, and identify every vertex $A \in V$ with an indicator vector $c_A \in \{0,1\}^d$ in the natural way.
Observe that for every $A,B \in V$ we have
\[M_{c_A, c_B} = f(c_A,c_B) = q(|A \cap B|).\]
Hence, for every $A \in V$ we have $|A|=s$ and thus $M_{c_A,c_A} =q(s) = 1$, whereas for every distinct non-adjacent $A,B \in V$ we have $m \leq |A \cap B|\leq s-1$ and thus $M_{c_A,c_B} = q(|A \cap B|) =0$.
Since the restriction of $M$ to $V \times V$ is symmetric and has rank at most $R$ over the reals, \expref{Lemma}{lemma:rankp_R} implies that ${\minrank}_{\Fset}(\Kneser{d}{s}{m}) \leq R$ for every finite field $\Fset$ and that the bound can be achieved by a symmetric matrix, as desired.
\end{proof}

\subsection{Applications}

We gather below several applications of \expref{Theorem}{thm:minrk_Kneser}.

\subsubsection{The Odd Alternating Cycle Conjecture over finite fields}

%
In this section we disprove %
\expref{Conjecture}{conj:alternating} over every finite field.
We will use the simple fact that
%
threshold Kneser graphs   %
do not contain short odd cycles, as stated below (see, e.\,g.,~\cite{Denley97,Haviv18free}).

\begin{lemma}\label{lemma:cycle_K}
Let $\ell \geq 3$ be an odd integer.
For every even integer $d$ and an integer $m \leq \frac{d}{2\ell}$, the graph $\Kneser{d}{\frac{d}{2}}{m}$ contains no odd cycle of length at most $\ell$.
\end{lemma}

We prove the following theorem, confirming \expref{Theorem}{thm:Intro_Cycles}.

\begin{theorem}\label{thm:Cycles}
For every odd integer $\ell \geq 3$ there exists $\delta = \delta(\ell) >0$ such that for every sufficiently large integer $n$, there exists an $n$-vertex graph $G$ with no odd cycle of length at most $\ell$ such that for every finite field $\Fset$,
\[{\minrank}_{\Fset}(G) \leq n^{1-\delta}.\]
Moreover, the bound on the minrank can be achieved by a symmetric matrix.
\end{theorem}

\begin{proof}
Fix an odd integer $\ell \geq 3$.
For an integer $d$ divisible by $2 \ell$, consider the graph $G = \Kneser{d}{\frac{d}{2}}{m}$ where $m = \frac{d}{2 \ell}$.
By \expref{Lemma}{lemma:cycle_K}, $G$ contains no odd cycle of length at most $\ell$.
As for the minrank parameter, \expref{Theorem}{thm:minrk_Kneser} implies that for every finite field $\Fset$,
\[{\minrank}_{\Fset}(G) \leq \sum_{i=0}^{d/2-m}\binom{d}{i} \leq 2^{H(\frac{1}{2}-\frac{m}{d}) \cdot d} = 2^{H(\frac{1}{2}-\frac{1}{2\ell}) \cdot d},\]
where $H$ stands for the binary entropy function.
Since $G$ has $|V| = \binom{d}{d/2} = 2^{(1-o(1)) \cdot d}$ vertices, for any $\delta>0$ such that $H(\frac{1}{2}-\frac{1}{2\ell}) < 1-\delta$ we have ${\minrank}_{\Fset}(G) \leq |V|^{1-\delta}$ for every sufficiently large integer $d$.
The proof is completed by considering, for every sufficiently large integer $n$, some $n$-vertex subgraph of the graph defined above, where $d$ is the smallest integer divisible by $2\ell$ such that $n \leq \binom{d}{d/2}$.
\end{proof}


\subsubsection{Triangle-free graphs and the orthogonality dimension over the binary field}\label{sec:OD_F_2_triangle}

%
Now we turn to the proof of  %
\expref{Theorem}{thm:IntroTriangleFree}.
Its proof adopts the following special case of a result due to Lempel~\cite{Lempel75}.
\begin{lemma}[\cite{Lempel75}]\label{lem:lempel}
Let $M$ %
be  %
an $n$ by $n$ symmetric matrix over the binary field $\Fset_2$ with at least one nonzero diagonal entry and rank $r$. Then, there exists an $n$ by $r$ matrix $B$ over $\Fset_2$ satisfying $M = B \cdot B^T$.
\end{lemma}

\begin{proof}[Proof of \expref{Theorem}{thm:IntroTriangleFree}]
Apply \expref{Theorem}{thm:Cycles} with $\ell = 3$ to obtain some $\delta >0$ such that for every sufficiently large integer $n$, there exist a triangle-free $n$-vertex graph $G$ and an $n$ by $n$ symmetric matrix $M$ over $\Fset_2$ of rank $r = {\rank}_{\Fset_2}(M) \leq n^{1-\delta}$ that represents $G$.
By \expref{Lemma}{lem:lempel}, there exists an $n$ by $r$ matrix $B$ over $\Fset_2$ satisfying $M = B \cdot B^T$.
By assigning the $i$-th row of $B$ to the $i$-th vertex of $G$ we get an $r$-dimensional orthogonal representation of $\overline{G}$ over $\Fset_2$, hence $\od(\overline{G},\Fset_2) \leq r \leq n^{1-\delta}$.
\end{proof}

\subsubsection{The vector chromatic number vs. minrank}

The vector chromatic number of graphs, introduced by Karger, Motwani, and Sudan in~\cite{KargerMS98}, is defined as follows.

\begin{definition}[Vector Chromatic Number]\label{def:chi_v}
For a graph $G=(V,E)$ the \emph{vector chromatic number} of $G$,
denoted by $\vchrom(G)$, is the minimal real value of $\kappa > 1$ such
that there exists an assignment of a unit vector $w_v$ to every
vertex $v \in V$ satisfying the inequality $\langle w_{v}, w_{v'} \rangle \leq
-\frac{1}{\kappa -1}$ whenever $v$ and $v'$ are adjacent in $G$.
\end{definition}

To prove \expref{Theorem}{thm:Intro_chi_v}, we need the following simple fact that relates the minrank of a graph to the minrank of its complement (see, e.\,g.,~\cite[Remark~2.2]{Peeters96}).

\begin{fact}\label{fact:minrank_comp}
For every field $\Fset$ and an $n$-vertex graph $G$, $\minrank_{\Fset}(G) \cdot \minrank_{\Fset}(\overline{G}) \geq n$.
\end{fact}

\begin{proof}[Proof of \expref{Theorem}{thm:Intro_chi_v}]
For an integer $d$ divisible by $8$, consider the graph $G = \Kneser{d}{\frac{d}{2}}{m}$ where $m = \frac{d}{8}$.
We first claim that $\vchrom(G) \leq 3$.
To see this, assign to every vertex $A$ of $G$, representing a $\frac{d}{2}$-subset of $[d]$, the unit vector $w_A \in \R^d$ defined by $(w_A)_i = \frac{1}{\sqrt{d}}$ if $i \in A$ and $(w_A)_i = - \frac{1}{\sqrt{d}}$ otherwise.
Observe that every two distinct vertices $A$ and $B$ that are adjacent in $G$ satisfy $|A \cap B| < \frac{d}{8}$ and thus $|A \bigtriangleup B| > \frac{3d}{4}$, implying that $\langle w_A, w_B \rangle = \frac{d-2 \cdot |A \bigtriangleup B|}{d} < -\frac{1}{2}$. This implies that $\vchrom(G) \leq 3$, as claimed.
As for the minrank parameter, \expref{Theorem}{thm:minrk_Kneser} implies that for every finite field $\Fset$,
\[{\minrank}_{\Fset}(G) \leq \sum_{i=0}^{d/2-m}\binom{d}{i} \leq 2^{H(\frac{1}{2}-\frac{m}{d}) \cdot d} = 2^{H(3/8) \cdot d},\]
where $H$ stands for the binary entropy function.
Since $G$ has $n = \binom{d}{d/2} = 2^{(1-o(1)) \cdot d}$ vertices, for any $\delta < 1- H(3/8)$ we have ${\minrank}_{\Fset}(G) \leq n^{1-\delta}$ assuming that $d$ is sufficiently large. By \expref{Fact}{fact:minrank_comp}, this implies that
${\minrank}_{\Fset}(\overline{G}) \geq n^{\delta}$, and we are done.
\end{proof}

\section*{Acknowledgements}
%
%
We would like to thank L{\'{a}}szl{\'{o}} Babai, Jop Bri{\"{e}}t, and the anonymous referees for their very helpful suggestions that improved the presentation of the paper.

\bibliographystyle{tocplain}   %

%
%
%
%
\bibliography{bibstrings,v018a022,bibtail}

%
%
%
%
%
\begin{tocauthors}
\begin{tocinfo}[golovnev]
 Alexander Golovnev\\
 Assistant professor\\
 Department of Computer Science\\
 Georgetown University\\
 Washington, D.\,C, USA\\
 alexgolovnev\tocat{}gamil\tocdot{}com\\
 \url{https://golovnev.org/}
\end{tocinfo}
\begin{tocinfo}[haviv]
  Ishay Haviv\\
  Associate professor\\
  The Academic College of Tel Aviv-Yaffo\\
  Tel Aviv, Israel\\
  ishayhav\tocat{}mta\tocdot{}ac\tocdot{}il \\
  \url{http://www2.mta.ac.il/~ishayhav}
\end{tocinfo}
\end{tocauthors}

%
%
%
%
%
%
%
%
%
%
%
%
%
%
%
%
%
%
%
%
\begin{tocaboutauthors}
\begin{tocabout}[golovnev]  %
  \textsc{Alexander Golovnev} is an assistant professor at Georgetown University.  Prior to this,  he was a Rabin postdoctoral fellow at Harvard University,  a postdoctoral fellow at Columbia University,  and a research scientist at Yahoo Research.  He received his Ph.\,D. from the Courant Institute of Mathematical Sciences at New York University where he was supervised by \href{http://cims.nyu.edu/~regev/}{Oded Regev} and \href{https://cs.nyu.edu/~dodis/}{Yevgeniy Dodis}.  His interest in computational lower bounds was instilled by \href{https://alexanderskulikov.github.io/}{Alexander S.~Kulikov}---largely against his will.
%
%
\end{tocabout}
\begin{tocabout}[haviv]
\textsc{Ishay Haviv} is an associate professor at the school of Computer Science in the Academic College of Tel Aviv-Yaffo.
He received his Ph.\,D. from Tel Aviv University where he was supervised by \href{http://cims.nyu.edu/~regev/}{Oded Regev}. His research lies in the theoretical aspects of Computer Science and in their interactions with other fields like Information Theory and %
%
classical areas of mathematics.  %
\end{tocabout}
\end{tocaboutauthors}

\end{document}

